\section{Zadanie 3 (Dominik Korsa)}
Dana jest moneta symetryczna. Zaproponuj algorytm, który będzie symulował
zachowanie monety asymetrycznej o ustalonym wymiernym \(p \in (0, 1)\).
Jaka jest oczekiwana liczba rzutów (symetryczną monetą) przed zwróceniem wyniku
(symulowanej asymetrycznej monety)? Im mniejsza oczekiwana liczba rzutów, tym lepiej.

\subsection*{Wprowadzenie do rozwiązania}
Niech \(p\) oznacza prawdopodobieństwo, że wynikiem rzutu symulowaną asymetryczną
monetą będzie orzeł.

Niech \(\frac{a}{b} = p\), gdzie \(a, b \in \mathbb{N}_1\), \(a < b\)
oraz \(a\) i \(b\) są względnie pierwsze.
Niech \(N = \lceil \log_2 b \rceil\).

Zakładam, że w kroku algorytmu mogę zwrócić jedynie ostatnio wyrzuconą wartość.
Jeżeli mogę zwrócić dowolną, to \(N+1\) rzut w zaproponowanym algorytmie nie
jest konieczny.

\subsection*{Algorytm}
W \(i\)-tym kroku (\(i \in \mathbb{N}_1\)):
\begin{itemize}
  \item Rzuć monetą \(N+1\) razy.
  \item Potraktuj pierwsze \(N\) wyrzuconych monet jako zapis binarny liczby \(x\)
  (traktując orła jako 1 i reszkę jako 0). \(x\) będzie z przedziału
  \(x \in [0, 2^N) \supseteq [0, b)\).
  Każda wartość z tego przedziału jest równie prawdopodobna.
  \item W zależności od \(x\):
  \begin{itemize}
    \item Zdarzenie \(A_i\) -- \(x \in [0, a)\) i ostatni rzut to orzeł: zwróć orła.
    \item Zdarzenie \(B_i\) -- \(x \in [a, b)\) i ostatni rzut to reszka: zwróć reszkę.
    \item Jeżeli nie zaszło $A_i$ ani $B_i$: powtórz rzuty (przejdź do kolejnego kroku).
  \end{itemize}
\end{itemize}

\subsection*{Analiza prawdopodobieństwa}
Niech \(C_i = A_i \cup B_i\).

Niech \(X\) to zmienna losowa oznaczająca numer (\(i\)) pierwszego kroku w którym zaszło
zdarzenie \(C_i\).

Zauważmy, że te losowania są \textbf{bez pamięci} -- wynik losowania \(i\)
nie zależy od wyników poprzednich losowań. Mamy więc:
\begin{itemize}
  \item \(
    P(A_i \mid X \geq i)
    = \underbrace{\frac{a}{2^N}}_{x \in [0, a)}
    \cdot \underbrace{\frac{1}{2}}_{\text{ostatni rzut to orzeł}}
    = \frac{a}{2^{N+1}}
  \)
  \item \(
    P(B_i \mid X \geq i)
    = \underbrace{\frac{b-a}{2^N}}_{x \in [a, b)}
    \cdot \underbrace{\frac{1}{2}}_{\text{ostatni rzut to reszka}}
    = \frac{b-a}{2^{N+1}}
  \)
  \item \(A_i\) i \(B_i\) są rozłączne więc \(
    P(C_i \mid X \geq i)
    = P(A_i \mid X \geq i) + P(B_i \mid X \geq i)
    = \frac{b}{2^{N+1}}
  \)
\end{itemize}

\(X\) ma rozkład geometryczny o parametrze \(q = \frac{b}{2^{N+1}}\):
\[
  P(X = n) = P\Bigg( (\Omega \setminus C_1) \cap \cdots \cap (\Omega \setminus C_{n-1}) \cap C_n \Bigg)
  = (1-q)^{n-1} \cdot q
\]
Z wykładu wiemy, że wartość oczekiwana \(X\) w tym przypadku to \(E[X] = \frac{1}{q} = \frac{2^{N+1}}{b}\).

Niech \(Y\) oznacza oczekiwaną liczbę rzutów monetą.
W każdym kroku rzucaliśmy monetą $N+1$ razy, więc:
\[
E[Y] = (N+1) \cdot E[X] = (N+1) \cdot \frac{2^{N+1}}{b}
\]

\subsection*{Odpowiedź}
Jeżeli \(p = \frac{a}{b}\) oraz \(N = \lceil \log_2 b \rceil\) to wynikiem jest:
\[
\boxed{(N+1) \cdot \frac{2^{N+1}}{b}}
\]

W wersji zadania w której w kroku algorytmu mogę zwrócić dowolną wartość,
nie tylko ostatnio wyrzucona, wynikiem jest:
\[
\boxed{N \cdot \frac{2^N}{b}}
\]
Dla tej wersji nie ma dowodu powyżej, ale można go łatwo uprościć.
